\documentclass[11pt]{article}
\usepackage[a4paper,margin=25mm]{geometry}
\usepackage{amsmath,amssymb,amsthm,booktabs,tabularx,longtable}
\usepackage[T1]{fontenc}
\usepackage{lmodern}
\usepackage[hidelinks]{hyperref}
\newtheorem{theorem}{Theorem}
\newtheorem{corollary}[theorem]{Corollary}
\DeclareMathOperator{\tr}{tr}
\DeclareMathOperator{\rank}{rank}
\newcommand{\eps}{\varepsilon}
\title{Rank-two extensions of Spezia's 2018 determinant result:\\
the two conjectures in OEIS A322844,\\
Comments (revision 68)}
\author{Haobo Ma\\\small ChronoAI Pte Ltd, Singapore
\and Wenlin Zhang\\\small National University of Singapore\\
\small ChronoAI Pte Ltd, Singapore}
\date{November 5, 2026}
\begin{document}
\maketitle
\begin{abstract}
Extending Spezia's determinant result (2018), we prove his two conjectures
in OEIS A322844, Comments, the two statements in the conjectures block
(revision 68, August 23, 2025): the absolute second exterior trace of
the serpentine matrix is the stated quartic quasipolynomial, and every
exterior trace of order greater than two vanishes.
The matrix has rank exactly two for every size at least two.
A factorization through a two-dimensional space gives its characteristic
polynomial, including the sign of the second exterior trace.
Corollaries give a minimal polynomial, a permanent formula, diagonal sums,
generating functions and reciprocal sums. We distinguish the exponential
generating function from a shifted-factorial expression in the source.
\end{abstract}

\section{The matrix and the two conjectures}
Spezia's two conjectures in A322844, Comments, revision 68~\cite{s322844},
are proved below by a rank-two factorization. The earlier determinant
formula in A322277 was explicitly recorded as proved on December 1,
2018~\cite{s322277}. It supplies the top exterior trace, but not the
intermediate orders $3\leq k<n$ or the second-trace formula. Our result
extends that determinant formula; the existing scalar and sum formulas
below retain their source attribution.

For $n\geq1$, number rows and columns from $1$ to $n$ and set
\begin{equation}\label{eq:matrix}
(M_n)_{ij}=\begin{cases}
j+n(i-1),&i\text{ odd},\\
ni-j+1,&i\text{ even}.
\end{cases}
\end{equation}
Thus the first row is $1,\ldots,n$, the second is $2n,\ldots,n+1$,
and the direction alternates. This is the matrix in A317614 and
A322277~\cite{s317614,s322277}.
Write $\eps=n\bmod2$ and define
\begin{equation}\label{eq:quantities}
T_n=\frac{n(n^2+\eps)}2,\qquad
a_n=\frac{n^2}{12}\bigl(3(1+n^2)-2(2+n^2)\eps\bigr)
=\begin{cases}
n^2(n^2+1)/4,&n\text{ even},\\
n^2(n^2-1)/12,&n\text{ odd}.
\end{cases}
\end{equation}
The values of $a_n$ start with $0,5,6,68,50,333$ for $n=1,\ldots,6$,
as in A322844~\cite{s322844}. Under the heading ``Conjectures: (Start)'',
the first assertion identifies $a(n)$ as ``the absolute value'' of the
second exterior trace; the second says the higher trace ``is equal to zero''.
More precisely, these are
$|\tr(\bigwedge^2 M_n)|=a_n$ for $n>1$ and
$\tr(\bigwedge^k M_n)=0$ for $k>2$, respectively.
The second-trace sequence was introduced on December 28, 2018;
the higher-trace assertion was added on January 6, 2019, in revision 5.
Revision 14 (January 31, 2019) grouped the two as conjectures.
For $k>n$ the exterior space is the zero vector space, whose trace is zero;
$\bigwedge^0 M_n$ is the identity on a one-dimensional space.

\begin{theorem}\label{thm:main}
For every integer $n\geq2$, the characteristic polynomial is
\begin{equation}\label{eq:charpoly}
\det(xI_n-M_n)=x^{n-2}(x^2-T_nx-a_n).
\end{equation}
The matrix $M_n$ has rank two and
\begin{equation}\label{eq:ext}
\tr(\textstyle\bigwedge^k M_n)=
\begin{cases}
1,&k=0,\\ T_n,&k=1,\\ -a_n,&k=2,\\ 0,&k\geq3.
\end{cases}
\end{equation}
In fact $\bigwedge^k M_n$ is the zero map for every $k\geq3$.
\end{theorem}

The two conjectures concern traces, rather than the trace of the ordinary
power $M_n^k$. Equation~\eqref{eq:ext} proves both, and fixes the sign
hidden by the absolute value in the first one.

\section{A two-dimensional factorization}
Put
\[
u_i=ni-\frac{n-1}{2},\qquad s_i=(-1)^{i+1},\qquad
w_i=i-\frac{n+1}{2}.
\]
Equation~\eqref{eq:matrix} becomes $(M_n)_{ij}=u_i+s_iw_j$.
Indeed, for odd $i$ this is $ni+j-n$, and for even $i$ it is $ni-j+1$.
Consequently, over $\mathbb Q$,
\begin{equation}\label{eq:factor}
M_n=UV,\qquad
U=\begin{pmatrix}u_1&s_1\\\vdots&\vdots\\u_n&s_n\end{pmatrix},
\qquad V=\begin{pmatrix}1&\cdots&1\\w_1&\cdots&w_n\end{pmatrix}.
\end{equation}
This already gives $\rank M_n\leq2$: every image vector is a linear
combination of $u$ and $s$. Any wedge of three or more such vectors is
zero, proving the last assertion of the theorem.

The needed sums are
\begin{align*}
\sum_i w_i&=0,&\sum_i w_i^2&=\frac{n(n^2-1)}{12},\\
\sum_i s_i&=\eps,&\sum_i w_is_i&=-\frac{n(1-\eps)}2.
\end{align*}
For completeness, the first two follow on substituting
$\sum i=n(n+1)/2$ and $\sum i^2=n(n+1)(2n+1)/6$ into the definition
of $w_i$. Pairing successive terms gives
$\sum i s_i=-n/2$ if $n$ is even and $(n+1)/2$ if $n$ is odd;
subtracting $(n+1)\sum s_i/2$ gives the last sum.
Since $u_i=nw_i+(n^2+1)/2$, these identities give
\begin{equation}\label{eq:small}
B:=VU=\begin{pmatrix}
 n(n^2+1)/2&\eps\\
 n^2(n^2-1)/12&-n(1-\eps)/2
\end{pmatrix}.
\end{equation}
In particular $\tr B=T_n$. Its determinant is
\[
-\frac{n^2(n^2+1)(1-\eps)}4
-\frac{\eps n^2(n^2-1)}{12}=-a_n.
\]
Both parity cases in~\eqref{eq:quantities} show $a_n>0$ when $n\geq2$.
Thus $B$ is invertible, so $U$ is injective and $V$ is surjective.
It follows that $\operatorname{im}(UV)=\operatorname{im}U$ has dimension two.

To obtain the characteristic polynomial without any spectral assumption,
evaluate the block determinant
\[
\det\begin{pmatrix}I_n&U\\tV&I_2\end{pmatrix}
=\det(I_n-tUV)=\det(I_2-tVU).
\]
Eliminating either off-diagonal block by a block triangular matrix of
determinant one proves both equalities over $\mathbb Q[t]$.
This is the rectangular determinant identity underlying the usual
Cauchy--Binet formulations~\cite{knill}.
Substitute $t=x^{-1}$ in $\mathbb Q(x)$ and multiply by $x^n$;
equation~\eqref{eq:charpoly} follows from~\eqref{eq:small}.

Finally, in the standard basis indexed by the $k$-subsets of
$\{1,\ldots,n\}$, the diagonal entries of $\bigwedge^k M_n$ are the
principal $k\times k$ minors of $M_n$.
Their sum is the coefficient of $t^k$ in $\det(I_n+tM_n)$:
in the determinant expansion choose a matrix entry in precisely those
$k$ rows and an identity entry in the other rows.
Thus this trace is $(-1)^k$ times the coefficient of $x^{n-k}$
in~\eqref{eq:charpoly}. This proves~\eqref{eq:ext} for $k\leq n$;
the zero-space convention covers $k>n$.

For $n=1$, $M_1=(1)$, so the exterior traces are $1,1,0,0,\ldots$.
For the empty matrix $M_0$, they are $1,0,0,\ldots$.
These boundary cases also satisfy the higher-order conjecture.

\begin{corollary}\label{cor:spectrum}
For $n\geq2$, the nonzero eigenvalues are
$\lambda_\pm=(T_n\pm\sqrt{T_n^2+4a_n})/2$, with
$\lambda_-<0<\lambda_+$; the remaining $n-2$ eigenvalues are zero.
The matrix is diagonalizable over $\mathbb R$.
Its determinant is $-5$ for $n=2$ and zero for $n\geq3$.
\end{corollary}
\begin{proof}
The quadratic in~\eqref{eq:charpoly} has two distinct real roots with
negative product. Rank two gives a nullspace of dimension $n-2$,
so the eigenspaces have total dimension $n$. The determinant follows from
the constant term, or from rank; at $n=2$, $a_2=5$.
The determinant-zero observation was already recorded by Spezia in
A322277 (2018)~\cite{s322277}.
\end{proof}

\begin{corollary}\label{cor:powers}
For $n=2$ the minimal polynomial of $M_n$ is $x^2-T_nx-a_n$;
for $n\geq3$ it is $x(x^2-T_nx-a_n)$. In particular,
\[
M_n^{r+2}=T_nM_n^{r+1}+a_nM_n^r\qquad(r\geq1).
\]
The traces of ordinary powers satisfy the same recurrence for $r\geq1$,
with $p_1=T_n$ and $p_2=T_n^2+2a_n$.
\end{corollary}
\begin{proof}
The two nonzero eigenvalues are distinct, and zero is an eigenvalue
exactly when $n\geq3$. The minimal polynomial of a diagonalizable matrix
is the product of its distinct eigenvalue factors. Multiplication by
$M_n^{r-1}$ and taking traces gives the assertions; the value of $p_2$
uses $\lambda_++\lambda_-=T_n$ and $\lambda_+\lambda_-=-a_n$.
\end{proof}

\begin{corollary}\label{cor:permanent}
Let $e_k(w)$ be the elementary symmetric polynomial of degree $k$ in
$w_1,\ldots,w_n$, with $e_0=1$, and let
$P_n(z)=\prod_{i=1}^n(u_i+s_i z)$. Then the permanent in A322277 is
\[
\operatorname{per}M_n=
\sum_{k=0}^n k!(n-k)!\,e_k(w)\,[z^k]P_n(z).
\]
Only even $k$ contribute.
\end{corollary}
\begin{proof}
Expand each permutation product using $M_{ij}=u_i+s_iw_j$.
Choose $k$ rows for the second summand. For each $k$-subset of columns,
the assignments contribute $k!(n-k)!$ times its $w$-product.
Summing the row choices gives $[z^k]P_n$, and the column choices give
$e_k(w)$. Since the list of $w_i$ is unchanged by negation, its odd
elementary symmetric functions vanish.
\end{proof}

\section{Sums in the linked entries}
The factorization also gives the following identities for $1\leq i,j\leq n$:
\begin{equation}\label{eq:rowcol}
\sum_j(M_n)_{ij}=n^2i-\frac{n(n-1)}2,\qquad
\sum_i(M_n)_{ij}=\frac{n(n^2+1)}2+
\eps\left(j-\frac{n+1}{2}\right).
\end{equation}
They follow by summing $u_i+s_iw_j$ and using the sums above.
These are the row and column formulas in A241016 and
A317617~\cite{s241016,s317617}; reversing a row does not change its sum.

For diagonal offset $0\leq k<n$, put $\ell=n-k$ and
$\delta=\ell\bmod2$. Define
$D^+_{n,k}=\sum_{i=1}^{\ell}(M_n)_{i,i+k}$ and
$D^-_{n,k}=\sum_{i=1}^{\ell}(M_n)_{i+k,i}$. Then
\begin{align}
D^+_{n,k}&=\frac n4\bigl(2\ell^2+1-(-1)^\ell\bigr),\label{eq:upper}\\
D^-_{n,k}&=\frac12\bigl(n\ell(n+k)+\ell(1-(-1)^k)
                   +(-1)^k(n-2k)\delta\bigr).\label{eq:lower}
\end{align}
To check these formulas, use
$\sum_{i=1}^\ell s_i=\delta$ and
$\sum_{i=1}^\ell i s_i=(-\ell+(2\ell+1)\delta)/2$.
For the upper diagonal, summation gives
$\ell(n\ell+1)/2+(-\ell+n\delta)/2
=n(\ell^2+\delta)/2$.
For the lower diagonal it gives
$\ell(n(\ell+2k)+1)/2+
(-1)^k(-\ell+(2\ell-n)\delta)/2$,
which is~\eqref{eq:lower}.
These recover A325516 and A325655~\cite{s325516,s325655}.
At $k=1$, equation~\eqref{eq:upper} is the formula in
A323724~\cite{s323724}. At $k=0$, both are $T_n$, the trace A317614.
With diagonal offsets starting at zero, the first-column identity in
A325516 therefore reads $D^+_{n,0}=T_n$.

A323723 is the case $k=1$ of~\eqref{eq:lower}~\cite{s323723}.
Both nearest-diagonal sums are even: for $n=2m$ they equal
$2m(2m^2-2m+1)$ and $4m^3$, and for $n=2m+1$ they equal
$2m^2(2m+1)$ and $2m(2m^2+3m+2)$, respectively.
For the upper sum divided by two, the parity formulas reduce modulo five
to $m(2m^2-2m+1)$ and $m^2(2m+1)$. For $m=0,1,2,3,4$ their
residue sets are $\{0,1,4\}$ and $\{0,3,4\}$, so neither has residue
two. Thus the upper sum never ends in four, as stated in A323724.

The two parity polynomials in~\eqref{eq:quantities} also explain the
links to the standard polynomial sequences A002415 and
A037270~\cite{s002415,s037270}:
$a_n$ is the A002415 value $n^2(n^2-1)/12$ at odd $n$;
at even $n$, $2a_n=n^2(n^2+1)/2$ is the normal magic-square total.
These equalities follow directly from the displayed polynomials, not
from a new assertion about those sequences.
The rank sequence, including the empty matrix, is $\min(n,2)$, as linked
from A317614 to A130130~\cite{s130130}. This identifies the sequence;
the other interpretations of A130130 are separate statements.

\begin{corollary}\label{cor:triangles}
The row sums $U_n=\sum_{k=0}^{n-1}D^+_{n,k}$ and
$L_n=\sum_{k=0}^{n-1}D^-_{n,k}$ are
\begin{align*}
U_n&=\frac{n^2(n+1)(2n+1)+3n(n+\eps)}{12},\\
L_n&=\frac{n^2(n^2+1)}2+T_n-U_n.
\end{align*}
These are the defining formulas of A325517 and A325656
(Spezia, 2019)~\cite{s325517,s325656}. Also
$D^+_{n,n-1}=n$, while $D^-_{n,n-1}$ equals $n^2$ for even $n$
and $n^2-n+1$ for odd $n$, as in A325657~\cite{s325657}.
Writing $Q(z)=(1-z)^5(1+z)^2$, one has
\begin{align*}
\sum_{n\geq0}U_nz^n&=\frac{z(1+3z+7z^2+3z^3+2z^4)}{Q(z)},\\
\sum_{n\geq0}L_nz^n&=\frac{z(1+5z+13z^2+9z^3+4z^4)}{Q(z)},\\
\sum_{n\geq0}L_n\frac{z^n}{n!}
&=\frac{z}{24}\left(3e^{-z}+(21+78z+54z^2+8z^3)e^z\right).
\end{align*}
\end{corollary}
\begin{proof}
Sum $D^+_{n,k}=n(\ell^2+(\ell\bmod2))/2$ over $1\leq\ell\leq n$,
using $\sum\ell^2=n(n+1)(2n+1)/6$ and
$\sum(\ell\bmod2)=(n+\eps)/2$. The total of all entries is
$n^2(n^2+1)/2$; adding $T_n$ counts the main diagonal twice and
gives $U_n+L_n$. The last diagonal has a single entry, read directly
from~\eqref{eq:matrix}. To obtain the generating functions, apply
$\Theta=z\,d/dz$ to $(1-z)^{-1}$ and $e^z$ for each polynomial
in $n$, and replace $z$ by $-z$ for its alternating part.
For example
$24L_n=8n^4+6n^3+4n^2+3n-3n(-1)^n$, which gives the last display.
\end{proof}
The displayed replacements for the two generating functions of A325656
use $13z^2$ and $8z^3$, respectively; the source prints the products
\texttt{13*x*2} and \texttt{8*exp(2*x)*x*3} instead.
At $n=2$ those literal expressions give $34$ and $10$, respectively,
rather than $L_2=8$ (using ordinary and exponential coefficients).

\section{Scalar formulas in A322844}
\begin{corollary}\label{cor:generating}
Both step-two recurrences in A322844 hold, and
\[
\sum_{n\geq0}a_nz^n=
\frac{5z^2+6z^3+43z^4+20z^5+43z^6+6z^7+5z^8}{(1-z^2)^5}.
\]
For $b_m=a_{2m}$ and $c_m=a_{2m+1}$, their ordinary generating functions are
\[
\sum b_mz^m=\frac{z(5+43z+43z^2+5z^3)}{(1-z)^5},\qquad
\sum c_mz^m=\frac{2z(3+10z+3z^2)}{(1-z)^5}.
\]
With the convention $F(z)=\sum_{n\geq0}a_nz^n/n!$, the exponential
generating functions are
\begin{align}
F(z)&=\frac{e^z(2z^4+12z^3+15z^2+3z)
+e^{-z}(z^4-6z^3+9z^2-3z)}{12},\label{eq:egf}\\
\sum b_m\frac{z^m}{m!}&=e^z(4z^4+24z^3+29z^2+5z),\label{eq:evenegf}\\
\sum c_m\frac{z^m}{m!}&=\frac{e^z}{3}(4z^4+32z^3+57z^2+18z).
\label{eq:oddegf}
\end{align}
\end{corollary}
\begin{proof}
The parity polynomials in~\eqref{eq:quantities} have degree four,
so their fifth differences vanish. Multiplying each ordinary series by
its displayed denominator then leaves only the listed initial coefficients.
For the exponential formulas, write
\[
a_n=\frac{2n^4+n^2+(-1)^n(n^4+2n^2)}{12},\quad
b_m=4m^4+m^2,\quad
c_m=\frac{4m^4+8m^3+5m^2+m}{3}.
\]
Apply $\Theta^r e^z=\sum n^r z^n/n!$, and evaluate the alternating
terms at $-z$. This gives all three formulas.
\end{proof}

The expression labelled ``E.g.f.'' in A322844 is, instead,
\begin{equation}\label{eq:shifted}
\begin{aligned}
H(z)&=\sum_{n\geq0}a_n\frac{z^n}{(n+2)!},\\
12z^2H(z)&=e^{-z}(24+21z+9z^2+2z^3+z^4)-60\\
&\quad+e^z(36-33z+15z^2-4z^3+2z^4).
\end{aligned}
\end{equation}
The singularity at zero is removable. To verify the equality, differentiate
$z^2H(z)$ twice: the result is \eqref{eq:egf}, and $z^2H(z)$ and its
first derivative vanish at zero. Thus the constant and linear integration
terms are fixed. In particular $[z^2]H=5/24$, not $a_2/2!=5/2$;
\eqref{eq:egf} is the corrected conventional exponential generating function.
The source's two hyperbolic expressions have the precise interpretations
\begin{align*}
\sum b_m\frac{z^m}{(2m)!}
&=\frac{z(8+z)\cosh\sqrt z+2\sqrt z(1+3z)\sinh\sqrt z}{4},\\
\sum c_m\frac{z^m}{(2m+1)!}
&=\frac{6z\cosh\sqrt z+\sqrt z(6+z)\sinh\sqrt z}{12}.
\end{align*}
Their power-series expansions contain only integral powers of $z$.
Expanding the hyperbolic functions proves these two identities
coefficientwise from $b_m$ and $c_m$; the conventional
subsequence exponential series are \eqref{eq:evenegf}--\eqref{eq:oddegf}.

\begin{corollary}\label{cor:reciprocal}
The three reciprocal sums in A322844 are
\begin{align*}
\sum_{k\geq1}\frac1{a_{2k}}&=2+\frac{\pi^2}{6}-\pi\coth(\pi/2),\\
\sum_{k\geq1}\frac1{a_{2k+1}}&=15-\frac{3\pi^2}{2},\\
\sum_{n\geq2}\frac1{a_n}&=17-\frac{4\pi^2}{3}-\pi\coth(\pi/2).
\end{align*}
\end{corollary}
\begin{proof}
Direct partial fractions give
\[
\frac1{a_{2k}}=\frac1{k^2}-\frac1{k^2+1/4},\qquad
\frac1{a_{2k+1}}=\frac3{k(k+1)}-\frac{12}{(2k+1)^2}.
\]
All series converge absolutely. The partial-fraction expansion
$\coth z=z^{-1}+2z\sum_{k\geq1}(z^2+\pi^2k^2)^{-1}$
in DLMF 4.36.3~\cite{dlmf} at $z=\pi/2$ yields
$\sum(k^2+1/4)^{-1}=\pi\coth(\pi/2)-2$.
At zero, comparison of the linear coefficient with
$\coth z=z^{-1}+z/3+O(z^3)$ gives $\sum k^{-2}=\pi^2/6$.
This coefficient comparison is valid because the series after removing
the pole converges uniformly on $|z|\leq\pi/2$.
Hence $\sum_{k\geq1}(2k+1)^{-2}=\pi^2/8-1$;
the telescoping sum $\sum_{k\geq1}1/(k(k+1))=1$ proves the second
formula. Adding proves the third.
\end{proof}

\section{Two linked rank-two families}
The same reduction proves the characteristic-polynomial formulas recorded
by Benoit Cloitre (November 9, 2002) and Michael Somos (November 14, 2002)
in A002415, Comments~\cite{s002415}.

\begin{corollary}\label{cor:linked}
Put $h_n=n^2(n^2-1)/12$ and $g_n=n^2(n^2+1)/2$. For $n\geq2$, let
$N^+_{ij}=i+j$ and $N^-_{ij}=i-j$, with $1\leq i,j\leq n$. Then
\[
\det(xI-N^+)=x^{n-2}(x^2-n(n+1)x-h_n),\qquad
\det(xI-N^-)=x^{n-2}(x^2+h_n).
\]
Both matrices have rank two, and
$h_n=\tfrac12\sum_{i,j=1}^n(i-j)^2$ (Bala, 2007, A002415).
The ordinary and exponential generating functions of $h_n$ are
\[
\sum_{n\geq0}h_nz^n=\frac{z^2(1+z)}{(1-z)^5},\qquad
\sum_{n\geq0}h_n\frac{z^n}{n!}
=\frac{e^zz^2(6+6z+z^2)}{12}.
\]
The reciprocal formulas recorded in A002415 and A037270 are
\begin{align*}
\sum_{n\geq2}\frac1{h_n}&=21-2\pi^2,&
\sum_{n\geq2}\frac{(-1)^n}{h_n}&=\pi^2-9,\\
\sum_{n\geq1}\frac1{g_n}&=1+\frac{\pi^2}{3}-\pi\coth\pi,&
\sum_{n\geq1}\frac{(-1)^{n+1}}{g_n}
&=\pi\operatorname{csch}\pi+\frac{\pi^2}{6}-1.
\end{align*}
The nonalternating identities are attributed to Kotesovec (2016 for $h_n$,
2019 for $g_n$), and the alternating ones to Eldar (2020 for $h_n$,
2021 for $g_n$), in those records.
\end{corollary}
\begin{proof}
Write $v=(1,\ldots,n)^T$ and let $\mathbf1$ be the all-ones vector.
The matrices are $v\mathbf1^T\pm\mathbf1v^T$.
Writing $S_1=\sum i=n(n+1)/2$, $S_2=\sum i^2=n(n+1)(2n+1)/6$,
the reduced products are
$\left(\begin{smallmatrix}S_1&n\\ \pm S_2&\pm S_1\end{smallmatrix}\right)$.
Their determinants are $\mp h_n$ since $nS_2-S_1^2=h_n>0$.
This proves the polynomial and rank assertions by the same block
determinant argument. Expanding the squared differences proves Bala's
identity. Applying $\Theta^4-\Theta^2$ to $(1-z)^{-1}$ and $e^z$
and dividing by 12 gives the two generating functions.

For reciprocal sums, use
$h_n^{-1}=12((n^2-1)^{-1}-n^{-2})$ for $n\geq2$ and
$g_n^{-1}=2(n^{-2}-(n^2+1)^{-1})$.
Partial fractions give $\sum_{n\geq2}(n^2-1)^{-1}=3/4$ and its
alternating counterpart $\sum_{n\geq2}(-1)^n(n^2-1)^{-1}=1/4$;
both follow by shifting $1/(2(n-1))-1/(2(n+1))$.
The coth expansion gives
$S(b):=\sum_{n\geq1}(n^2+b^2)^{-1}
=(\pi b\coth(\pi b)-1)/(2b^2)$.
Splitting even and odd terms gives the alternating value
$S(1)-S(1/2)/2=(1-\pi\operatorname{csch}\pi)/2$,
using $\coth(z/2)=\coth z+\operatorname{csch}z$.
Together with $\sum(-1)^{n+1}/n^2=\pi^2/12$, these prove all four
identities. Every rearrangement is justified by absolute convergence.
\end{proof}

Two indexing or identification points in the linked sources merit precision.
In A325516, Formula, ``1st column: T(n, 1) = A317614(n)''
uses $k=1$ where $k=0$ is needed:
already $D^+_{2,1}=2\ne T_2=4$.
The Crossrefs section of A317617 calls A074147 ``the elements of the
antidiagonal'' of $M$. This identification fails at size four.
Direct substitution into~\eqref{eq:matrix} gives
\begin{equation}\label{eq:anti}
(M_n)_{i,n+1-i}=\begin{cases}(n-1)i+1,&i\text{ odd},\\
(n+1)i-n,&i\text{ even}.
\end{cases}
\end{equation}
For $n=4$ these entries are $(4,6,10,16)$, whereas row four of the
triangle in A074147 is $(6,8,10,12)$~\cite{s074147}.
Thus the latter sequence is not the antidiagonal triangle of these matrices.

\section{Coverage and verification}
``Proved'' below applies to an assertion actually made by the cited source.
Consequences not asserted there are labelled ``not posed; proved''.
Dates denote the cited records, not a claim of priority.
\begingroup\small
\begin{longtable}{@{}p{0.30\linewidth}p{0.66\linewidth}@{}}\toprule
Source item & Mathematical coverage\\\midrule
A322844, Comments, two conjectures (rev.\ 68) & Proved: both quoted conjectures;
signed second trace and all higher traces
(Theorem~\ref{thm:main}).\\
A317614, trace & Reproved; prior formula Spezia (2018): $T_n$.\\
A322277, determinant & Reproved; prior result Spezia (December 1, 2018).
It already gives the top exterior trace for $n>2$.\\
A322277, permanent & Not posed; proved: coefficient formula
(Corollary~\ref{cor:permanent}); the source defines the permanent sequence.\\
A241016; A317617 & Reproved; prior formulas Ngaokrajang (2014),
Lang (2014), and Spezia (2018): \eqref{eq:rowcol}.\\
A325516; A323724 & Reproved; prior formulas Spezia (2019) and Berselli (2019).
A325516 first-column index refuted-and-corrected to $k=0$.\\
A325655; A323723 & Reproved; prior formulas Spezia (2019),
\eqref{eq:lower}.\\
A325517; A325656; A325657 & Reproved; prior formulas Spezia (2019):
diagonal aggregates and last entries (Corollary~\ref{cor:triangles}).
A325656 literal generating-function products refuted-and-corrected.\\
A130130 & Not posed there; proved: serpentine rank is this sequence,
$\min(n,2)$. No assertion about its other interpretations.\\
A002415; A037270 & Polynomial specializations identified above.
Reproved; prior matrix formulas Cloitre and Somos (2002), squared-difference
formula Bala (2007), and reciprocal formulas Kotesovec (2016, 2019),
Eldar (2020, 2021); Corollary~\ref{cor:linked}.\\
A317617 link to A074147 & Refuted-and-corrected: direct entry formula
\eqref{eq:anti} and a size-four counterexample.\\
A322844, Formula & Proved: parity formulas, recurrences, all three ordinary
generating functions and reciprocal sums
(Corollaries~\ref{cor:generating}, \ref{cor:reciprocal}).\\
A322844, single-sequence E.g.f. & Refuted-and-corrected for the conventional
$n!$ weights; exact shifted weights in \eqref{eq:shifted}.\\
A322844, parity hyperbolic expressions & Proved with the displayed
$(2m)!$ and $(2m+1)!$ weights; conventional subsequence series in
\eqref{eq:evenegf}--\eqref{eq:oddegf}.\\
Additional spectral claims & Not posed; proved: zero exterior maps,
diagonalizability, minimal polynomial and ordinary-power recurrence.\\
Other interpretations and bivariate series in linked entries & Not covered:
Moessner-type algorithms, unrelated sequence interpretations and bivariate
generating functions outside the triangular matrix-sum domain.\\\bottomrule
\end{longtable}
\endgroup

The arguments above are written proofs for all stated sizes.
The file \texttt{verify.py} additionally checks every principal minor
for $0\leq n\leq10$, the complete characteristic polynomial for
$1\leq n\leq16$, and the factorization, rank and diagonal-sum formulas
and antidiagonal entry formulas through $n=32$, using exact rational arithmetic.
Its characteristic-polynomial calculation uses traces of ordinary matrix
powers and Newton's recurrence, separately from the two-dimensional formula.
It also checks the permanent formula through size 9, the matrix-power
identity through size 16, all displayed scalar and aggregate generating
functions through degree 64, and exact reciprocal partial fractions for
$1\leq k\leq64$. A bounded numerical sum supplies a separate check of
the reciprocal constants; their all-size justification is the written proof.
The two linked matrix families' full characteristic polynomials, ranks
and squared-difference formula are checked through size 16.

Two pinned Lean declarations express the exterior-trace results using
characteristic coefficients. In
\href{https://github.com/the-omega-institute/trureturing/blob/d9bb8a406783c5222849b798f3c2ab738f653d75/D5/S3/Quantum/Matrix/SpeziaSerpentineSecondExteriorTrace.lean}
{\texttt{SpeziaSerpentineSecondExteriorTrace.lean}},
\texttt{result} states the absolute second-trace identity for $n>1$,
and \texttt{charpoly\_factor} gives the rectangular factorization of the
characteristic polynomial for $n\geq2$.
Its declaration \texttt{coefficient\_two} gives the signed coefficient
$-a_n$ for $n\geq2$.
In
\href{https://github.com/the-omega-institute/trureturing/blob/d9bb8a406783c5222849b798f3c2ab738f653d75/D5/S3/Quantum/Matrix/SpeziaSerpentineHigherExteriorTrace.lean}
{\texttt{SpeziaSerpentineHigherExteriorTrace.lean}},
\texttt{result} states that every exterior trace of order $k>2$ is zero,
including $k>n$. Both use the integer matrix~\eqref{eq:matrix}, with
Lean indices shifted by one. The zero-map assertion, diagonalizability
and row, column and diagonal identities have their written proofs above;
they are distinct from those two exterior-trace declarations.

\begingroup\footnotesize
\begin{thebibliography}{18}
\raggedright
\setlength{\itemsep}{0pt}
\bibitem{s322844} S. Spezia, OEIS A322844, December 28, 2018,
Comments (two conjectures), revision 68 (August 23, 2025),
\url{https://oeis.org/A322844}.
History: \url{https://oeis.org/history?seq=A322844}, revisions 2, 5 and 14.
\bibitem{s317614} S. Spezia, OEIS A317614, August 1, 2018,
\url{https://oeis.org/A317614}.
\bibitem{s322277} S. Spezia, OEIS A322277, December 1, 2018,
\url{https://oeis.org/A322277}.
\bibitem{s241016} K. Ngaokrajang, OEIS A241016, August 8, 2014,
\url{https://oeis.org/A241016}.
\bibitem{s317617} S. Spezia, OEIS A317617, August 1, 2018,
\url{https://oeis.org/A317617}.
\bibitem{s325516} S. Spezia, OEIS A325516, May 7, 2019,
\url{https://oeis.org/A325516}.
\bibitem{s325655} S. Spezia, OEIS A325655, May 13, 2019,
\url{https://oeis.org/A325655}.
\bibitem{s323723} S. Spezia, OEIS A323723, January 25, 2019,
\url{https://oeis.org/A323723}.
\bibitem{s323724} S. Spezia, OEIS A323724, January 25, 2019,
\url{https://oeis.org/A323724}.
\bibitem{s325517} S. Spezia, OEIS A325517, May 7, 2019,
\url{https://oeis.org/A325517}.
\bibitem{s325656} S. Spezia, OEIS A325656, May 13, 2019,
\url{https://oeis.org/A325656}.
\bibitem{s325657} S. Spezia, OEIS A325657, May 13, 2019,
\url{https://oeis.org/A325657}.
\bibitem{s002415} N. J. A. Sloane, OEIS A002415,
\url{https://oeis.org/A002415}.
\bibitem{s037270} A. Gulliver, OEIS A037270,
\url{https://oeis.org/A037270}.
\bibitem{s074147} A. Murthy, OEIS A074147, August 28, 2002,
\url{https://oeis.org/A074147}.
\bibitem{s130130} P. Curtz, OEIS A130130, August 1, 2007,
\url{https://oeis.org/A130130}.
\bibitem{knill} O. Knill, Cauchy--Binet for pseudo-determinants,
\emph{Linear Algebra Appl.} \textbf{459} (2014), 522--547;
arXiv:1306.0062v3, \url{https://arxiv.org/abs/1306.0062}.
\bibitem{dlmf} NIST Digital Library of Mathematical Functions,
Section 4.36, equation 4.36.3,
\url{https://dlmf.nist.gov/4.36.E3}.
\end{thebibliography}
\endgroup
\noindent\small\copyright\ 2026 Haobo Ma and Wenlin Zhang.
Licensed under Creative Commons Attribution 4.0 International
(CC BY 4.0), \url{https://creativecommons.org/licenses/by/4.0/}.
\end{document}
